% GENERATED FILE. Edit canonical lessons, then run npm run generate:books.
% canonical-source-hash: fnv1a64:7a61531e66ab4c97
% canonical-lessons: SA-C214-eva, SA-C214-khalu, SA-C214-iva, SA-C214-vina, SA-C214-prati

\chapter{Exactly, Indeed, Like, Without, Towards}
\label{ch:sa-exactly-indeed-like-without-towards}
\hlchaptermodality{214}

\begin{chapteropening}
\textbf{By the end of this chapter:} \emph{I can say exactly, indeed, like, without and towards.}

Complete the last lesson of chapter 214: I can say exactly, indeed, like, without and towards.
\end{chapteropening}

\section[\texorpdfstring{\sk{एव} (eva)}{एव (eva)}]{\sk{एव} (eva) — exactly, only}

\label{lesson:SA-C214-eva}



\begin{quote}
Before the new one: say the Sanskrit for afterwards, then the Sanskrit for as long as.
\end{quote}

\subsection*{You'll want to know: \sk{एव}}
\textbf{\sk{एव}} — \emph{eva} — \textquotedblleft{}exactly, only\textquotedblright{}.

\begin{grammarlens}[title={What you've built}]
One more word for this chapter.
\end{grammarlens}

\subsection*{Your turn}
\begin{itemize}
\raggedright
  \item \emph{Say it:} \emph{eva}
  \item \emph{Say it:} \emph{eva}, once more
  \item \emph{Say it:} say \emph{eva}
  \item \emph{Recall:} say the Sanskrit for afterwards, then the Sanskrit for as long as, then say \emph{eva} again
\end{itemize}

\begin{tcolorbox}[breakable,colback=teal!4,colframe=teal!35!black,title={Before you move on}]
What does \sk{एव} mean? (\textquotedblleft{}Exactly, only\textquotedblright{}.) Say it once more.
\end{tcolorbox}
\section[\texorpdfstring{\sk{खलु} (khalu)}{खलु (khalu)}]{\sk{खलु} (khalu) — indeed}

\label{lesson:SA-C214-khalu}



\begin{quote}
Before the new one: say the Sanskrit for as long as, then the Sanskrit for exactly.
\end{quote}

\subsection*{You'll want to know: \sk{खलु}}
\textbf{\sk{खलु}} — \emph{khalu} — \textquotedblleft{}indeed\textquotedblright{}.

\begin{grammarlens}[title={What you've built}]
One more word for this chapter.
\end{grammarlens}

\subsection*{Your turn}
\begin{itemize}
\raggedright
  \item \emph{Say it:} \emph{khalu}
  \item \emph{Say it:} \emph{khalu}, once more
  \item \emph{Say it:} say \emph{khalu}
  \item \emph{Recall:} say the Sanskrit for as long as, then the Sanskrit for exactly, then say \emph{khalu} again
\end{itemize}

\begin{tcolorbox}[breakable,colback=teal!4,colframe=teal!35!black,title={Before you move on}]
What does \sk{खलु} mean? (\textquotedblleft{}Indeed\textquotedblright{}.) Say it once more.
\end{tcolorbox}
\section[\texorpdfstring{\sk{इव} (iva)}{इव (iva)}]{\sk{इव} (iva) — like, as if}

\label{lesson:SA-C214-iva}



\begin{quote}
Before the new one: say the Sanskrit for exactly, then the Sanskrit for indeed.
\end{quote}

\subsection*{You'll want to know: \sk{इव}}
\textbf{\sk{इव}} — \emph{iva} — \textquotedblleft{}like, as if\textquotedblright{}.

\begin{grammarlens}[title={What you've built}]
One more word for this chapter.
\end{grammarlens}

\subsection*{Your turn}
\begin{itemize}
\raggedright
  \item \emph{Say it:} \emph{iva}
  \item \emph{Say it:} \emph{iva}, once more
  \item \emph{Say it:} say \emph{iva}
  \item \emph{Recall:} say the Sanskrit for exactly, then the Sanskrit for indeed, then say \emph{iva} again
\end{itemize}

\begin{tcolorbox}[breakable,colback=teal!4,colframe=teal!35!black,title={Before you move on}]
What does \sk{इव} mean? (\textquotedblleft{}Like, as if\textquotedblright{}.) Say it once more.
\end{tcolorbox}
\section[\texorpdfstring{\sk{विना} (vinā)}{विना (vinā)}]{\sk{विना} (vinā) — without}

\label{lesson:SA-C214-vina}



\begin{quote}
Before the new one: say the Sanskrit for indeed, then the Sanskrit for like.
\end{quote}

\subsection*{You'll want to know: \sk{विना}}
\textbf{\sk{विना}} — \emph{vinā} — \textquotedblleft{}without\textquotedblright{}.

\begin{grammarlens}[title={What you've built}]
One more word for this chapter.
\end{grammarlens}

\subsection*{Your turn}
\begin{itemize}
\raggedright
  \item \emph{Say it:} \emph{vinā}
  \item \emph{Say it:} \emph{vinā}, once more
  \item \emph{Say it:} say \emph{vinā}
  \item \emph{Recall:} say the Sanskrit for indeed, then the Sanskrit for like, then say \emph{vinā} again
\end{itemize}

\begin{tcolorbox}[breakable,colback=teal!4,colframe=teal!35!black,title={Before you move on}]
What does \sk{विना} mean? (\textquotedblleft{}Without\textquotedblright{}.) Say it once more.
\end{tcolorbox}
\section[\texorpdfstring{\sk{प्रति} (prati)}{प्रति (prati)}]{\sk{प्रति} (prati) — towards}

\label{lesson:SA-C214-prati}



\begin{quote}
Before the new one: say the Sanskrit for like, then the Sanskrit for without.
\end{quote}

\subsection*{You'll want to know: \sk{प्रति}}
\textbf{\sk{प्रति}} — \emph{prati} — \textquotedblleft{}towards\textquotedblright{}.

\begin{grammarlens}[title={What you've built}]
One more word for this chapter.
\end{grammarlens}

\subsection*{Your turn}
\begin{itemize}
\raggedright
  \item \emph{Say it:} \emph{prati}
  \item \emph{Say it:} \emph{prati}, once more
  \item \emph{Say it:} say \emph{prati}
  \item \emph{Recall:} say the Sanskrit for like, then the Sanskrit for without, then say \emph{prati} again
\end{itemize}

\begin{tcolorbox}[breakable,colback=teal!4,colframe=teal!35!black,title={Before you move on}]
What does \sk{प्रति} mean? (\textquotedblleft{}Towards\textquotedblright{}.) Say it once more.
\end{tcolorbox}
